%============================================================
% CH1 · P3 · PROBLEM 2 — MAX AREA OF ISLAND (LeetCode 695)
%============================================================
\subsection{Problem 1.3.2: Max Area of Island}
\problemlabel{1.3.2}{ssec:ch1-max-area}

\problemcard{1.3.2}{Max Area of Island}{LeetCode 695}{https://leetcode.com/problems/max-area-of-island/}{\exmedium\ \lc{Medium}}{Flood fill that returns the size of the component; maximum over all components}

\subsubsection*{1. Problem statement}

You are given an $R \times C$ binary matrix \texttt{grid}. An \textbf{island} is a group of \texttt{1}s (land) connected 4-directionally (horizontally or vertically); everything outside the grid is water. The \textbf{area} of an island is the number of its cells. Return the maximum area of an island, or $0$ if there is no island.

\textbf{Constraints:} $1 \le R, C \le 50$, \; \texttt{grid[i][j]} is $0$ or $1$.

\begin{center}
\small
\begin{tabular}{@{}lcl@{}}
\toprule
\textbf{grid} & \textbf{Output} & \textbf{Note} \\
\midrule
the $8\times13$ map below & 6 & not 11: two islands touch only at a corner \\
\texttt{[[0,0,0,0,0,0,0,0]]} & 0 & no land at all \\
\bottomrule
\end{tabular}
\end{center}

\begin{diagrambox}
\begin{tikzpicture}
    \def\gIpIIIs{0.55}
    \gIpIIIgrid{wwowwwwowwwww,wwwwwwwooowww,woowowwwwwwww,wowwoowwowoww,wowwoowwoooww,wwwwwwwwwwoww,wwwwwwwooowww,wwwwwwwoowwww}
    \foreach \c in {0,...,12} \node[font=\tiny, text=darkgray] at (0.55*\c, 0.42) {\c};
    \foreach \r in {0,...,7} \node[font=\tiny, text=darkgray] at (-0.45, -0.55*\r) {\r};
    \node[font=\small\bfseries] at (3.3,-4.5) {Example 1 (white = land, grey = water)};
\end{tikzpicture}
\figcap{The first example of the statement. How many islands do you see, and which is the largest?}
\end{diagrambox}

\subsubsection*{2. The graph model}

\begin{modelbox}[Graph Model --- Max Area of Island]
\renewcommand{\arraystretch}{1.35}
\begin{tabularx}{\linewidth}{@{}>{\bfseries}l X@{}}
Vertices & the land cells; at most $RC = 2500$ \\
Edges & between land cells sharing a side; undirected; implicit \\
Weights & none (every vertex counts 1 towards the area) \\
Question & the \textbf{largest number of vertices in a connected component} (0 if there is no vertex) \\
Size & $|V| \le 2500$, $|E| \le 5000$: anything up to $\bigO{(RC)^2} \approx 6\cdot10^6$ fits; $\bigO{RC}$ is the clean target \\
\end{tabularx}
\end{modelbox}

\begin{diagrambox}
\begin{tikzpicture}
    \def\gIpIIIs{0.55}
    \gIpIIIgrid{wwAwwwwBwwwww,wwwwwwwBBBwww,wCCwDwwwwwwww,wCwwDDwwEwEww,wCwwDDwwEEEww,wwwwwwwwwwEww,wwwwwwwFFFwww,wwwwwwwFFwwww}
    \draw[warningred, line width=1.4pt] (5.225,-3.025) circle (0.5);
    \draw[warningred, -{Stealth[length=2mm]}] (6.6,-3.6) -- (5.7,-3.2);
    \node[font=\scriptsize, text=warningred, anchor=west] at (6.6,-3.7) {corner contact: no edge};
    \node[font=\scriptsize, anchor=west] at (7.6,-1.2) {
        \begin{tabular}{@{}clc@{}}
        \toprule
        island & first cell & area \\
        \midrule
        A & $(0,2)$ & 1 \\
        B & $(0,7)$ & 4 \\
        C & $(2,1)$ & 4 \\
        D & $(2,4)$ & 5 \\
        E & $(3,8)$ & \textbf{6} \\
        F & $(6,7)$ & 5 \\
        \bottomrule
        \end{tabular}};
\end{tikzpicture}
\figcap{The components, named in the order the row-major scan meets them. E's cell $(5,10)$ and F's cell $(6,9)$ touch only at a corner, so they are \emph{different} islands: counting them together would give the wrong answer $6+5=11$.}
\end{diagrambox}

\subsubsection*{3. Key observations}

\begin{itemize}[leftmargin=*]
    \item \textbf{Area = size of a component.} Problem 1.3.1 counted components; here we \emph{measure} each one and keep the maximum. The only change: the fill returns how many cells it marked.
    \item \textbf{A recursive fill can return the size directly:} area $= 1 + $ the areas returned by the calls on the four neighbours. Each cell contributes its 1 exactly once because it is marked before the recursion continues.
    \item \textbf{Recursion is safe here:} at most $RC = 2500$ land cells, so the depth is at most 2500 frames --- tiny.
    \item \textbf{Sinking:} setting a visited land cell to $0$ removes it from the graph; no extra \texttt{visited} array is needed (LeetCode lets us modify \texttt{grid}).
\end{itemize}

\subsubsection*{4. Approaches}

\begin{bruteforcebox}[Approach 1: Brute Force --- a Fresh Fill From Every Land Cell]
\textbf{Idea.} For every land cell, start a new fill with a \emph{fresh} \texttt{seen} array and count the cells it reaches; the answer is the largest count. Every island is measured once per cell it contains.

\begin{lstlisting}[style=cpp]
#include <bits/stdc++.h>
using namespace std;

class Solution {
public:
    int maxAreaOfIsland(vector<vector<int>>& grid) {
        int R = grid.size(), C = grid[0].size();
        const int dr[4] = {-1, 1, 0, 0}, dc[4] = {0, 0, -1, 1};
        int best = 0;
        for (int r = 0; r < R; r++)
            for (int c = 0; c < C; c++) {
                if (grid[r][c] != 1) continue;
                vector<vector<char>> seen(R, vector<char>(C, 0));   // fresh every time
                vector<pair<int,int>> st(1, make_pair(r, c));
                seen[r][c] = 1;
                int area = 0;
                while (!st.empty()) {
                    int x = st.back().first, y = st.back().second;
                    st.pop_back();
                    area++;
                    for (int d = 0; d < 4; d++) {
                        int nx = x + dr[d], ny = y + dc[d];
                        if (nx < 0 || nx >= R || ny < 0 || ny >= C) continue;
                        if (grid[nx][ny] != 1 || seen[nx][ny]) continue;
                        seen[nx][ny] = 1;
                        st.push_back(make_pair(nx, ny));
                    }
                }
                best = max(best, area);
            }
        return best;
    }
};
\end{lstlisting}

\begin{diagrambox}
\begin{tikzpicture}
    \def\gIpIIIs{0.5}
    \foreach \x in {0, 2.6, 5.2} {
        \begin{scope}[xshift=\x cm]
            \gIpIIIgrid{aaw,aww,wwb}
        \end{scope}
    }
    \gIpIIIat{0}{0}{s}
    \begin{scope}[xshift=2.6cm]\gIpIIIat{0}{1}{s}\end{scope}
    \begin{scope}[xshift=5.2cm]\gIpIIIat{1}{0}{s}\end{scope}
    \node[note, anchor=west, text width=4.2cm, align=left] at (7.0,-0.5) {3 starts $\times$ 3 cells = 9 steps for the blue island, then one more fill for the green one.};
\end{tikzpicture}
\figcap{The same 3-cell island (blue) is filled again from each of its cells (s = start). An island of $k$ cells costs $k^2$ steps, plus a fresh $R\times C$ array per start.}
\end{diagrambox}

\begin{complexitybox}
\textbf{Time:} $\bigO{(RC)^2}$ --- each of up to $RC$ starts allocates and fills up to $RC$ cells. For $50\times50$ that is about $6\cdot10^6$: accepted on LeetCode, but the same idea on a $1000\times1000$ map would need $10^{12}$ steps. \textbf{Space:} $\bigO{RC}$. \textbf{Verdict:} \ac{} here, \tle{} in general --- and our reference for random tests.
\end{complexitybox}
\end{bruteforcebox}

\begin{algorithmbox}[Approach 2: Scan + Recursive Fill That Returns the Area]
\textbf{Idea.} Scan the grid; each land cell still present starts a fill that \emph{sinks} its island (sets every cell to $0$) and returns the number of cells sunk. Keep the maximum.

\begin{lstlisting}[style=cpp]
#include <bits/stdc++.h>
using namespace std;

class Solution {
    int R, C;
    int sink(vector<vector<int>>& g, int r, int c) {   // returns the area it sinks
        if (r < 0 || r >= R || c < 0 || c >= C || g[r][c] != 1) return 0;
        g[r][c] = 0;                                   // visited = turned into water
        return 1 + sink(g, r - 1, c) + sink(g, r + 1, c)
                 + sink(g, r, c - 1) + sink(g, r, c + 1);
    }
public:
    int maxAreaOfIsland(vector<vector<int>>& grid) {
        R = grid.size(); C = grid[0].size();
        int best = 0;
        for (int r = 0; r < R; r++)
            for (int c = 0; c < C; c++)
                if (grid[r][c] == 1) best = max(best, sink(grid, r, c));
        return best;
    }
};
\end{lstlisting}

Here the checks are made \emph{at the top of the call} (``am I outside, or not land?'') instead of before the call. That costs a few extra calls that return 0 immediately, but the fill stays three lines long. It is the standard style for small grids.

\begin{complexitybox}
\textbf{Time:} $\bigO{RC}$ --- every land cell is sunk once and makes 4 calls; every water cell is only looked at. \textbf{Space:} $\bigO{RC}$ recursion depth in the worst case ($\le 2500$ frames, safe). \textbf{Verdict:} \ac{}.
\end{complexitybox}
\end{algorithmbox}

\begin{invariantbox}[Why it is correct]
\textbf{Lemma.} If $(r,c)$ is land when \texttt{sink(r,c)} is called, the call sinks exactly the land cells that are reachable from $(r,c)$ through land \emph{at the moment of the call}, and returns their number.

\emph{Proof by induction on the number of land cells.} The call sinks $(r,c)$ itself (1 cell) and then calls its four neighbours one after another. Each call on a neighbour that is still land sinks, by induction (the grid now has fewer land cells), the part of the island reachable from that neighbour without passing through cells already sunk, and returns its size; calls on water return 0. The parts are disjoint, because a cell is sunk at most once, and together with $(r,c)$ they cover everything reachable from $(r,c)$: such a cell is reached through some neighbour of $(r,c)$, and is either sunk by that neighbour's call or already sunk by an earlier sibling call. So $1 + \sum$ counts every island cell exactly once.

\textbf{Main loop.} By the lemma, each call from the scan sinks one complete island and returns its area; islands sunk earlier are never measured again, and every island is met by the scan. Hence \texttt{best} is the maximum area. \hfill$\blacksquare$
\end{invariantbox}

\begin{alternativebox}[Approach 3: Alternative Algorithms --- BFS, or DSU With Sizes]
\textbf{BFS, input untouched.} If the grid must not be modified, use a separate \texttt{seen} array and count a cell when it leaves the queue:
\begin{lstlisting}[style=cpp]
#include <bits/stdc++.h>
using namespace std;

class Solution {
public:
    int maxAreaOfIsland(vector<vector<int>>& grid) {
        int R = grid.size(), C = grid[0].size();
        const int dr[4] = {-1, 1, 0, 0}, dc[4] = {0, 0, -1, 1};
        vector<vector<char>> seen(R, vector<char>(C, 0));
        int best = 0;
        for (int r = 0; r < R; r++)
            for (int c = 0; c < C; c++) {
                if (grid[r][c] != 1 || seen[r][c]) continue;
                queue<int> q;
                q.push(r * C + c);
                seen[r][c] = 1;
                int area = 0;
                while (!q.empty()) {
                    int id = q.front(); q.pop();
                    area++;                                    // count when DEQUEUED
                    int x = id / C, y = id % C;
                    for (int d = 0; d < 4; d++) {
                        int nx = x + dr[d], ny = y + dc[d];
                        if (nx < 0 || nx >= R || ny < 0 || ny >= C) continue;
                        if (grid[nx][ny] != 1 || seen[nx][ny]) continue;
                        seen[nx][ny] = 1;
                        q.push(nx * C + ny);
                    }
                }
                best = max(best, area);
            }
        return best;
    }
};
\end{lstlisting}

\textbf{DSU with sizes (preview of Chapter 6).} Unite every land cell with its land neighbours below and to the right; union by size keeps \texttt{sz[root]} = area of the island. The answer is the largest \texttt{sz} of a root.
\begin{lstlisting}[style=cpp]
#include <bits/stdc++.h>
using namespace std;

class Solution {
    vector<int> par, sz;
    int find(int x) { while (par[x] != x) x = par[x] = par[par[x]]; return x; }
    void unite(int a, int b) {
        a = find(a); b = find(b);
        if (a == b) return;
        if (sz[a] < sz[b]) swap(a, b);
        par[b] = a; sz[a] += sz[b];                    // union by size
    }
public:
    int maxAreaOfIsland(vector<vector<int>>& grid) {
        int R = grid.size(), C = grid[0].size();
        par.resize(R * C); sz.assign(R * C, 1);
        for (int i = 0; i < R * C; i++) par[i] = i;
        for (int r = 0; r < R; r++)
            for (int c = 0; c < C; c++) {
                if (grid[r][c] != 1) continue;
                if (r + 1 < R && grid[r + 1][c] == 1) unite(r * C + c, (r + 1) * C + c);
                if (c + 1 < C && grid[r][c + 1] == 1) unite(r * C + c, r * C + c + 1);
            }
        int best = 0;
        for (int i = 0; i < R * C; i++)
            if (grid[i / C][i % C] == 1 && find(i) == i) best = max(best, sz[i]);
        return best;
    }
};
\end{lstlisting}

\begin{diagrambox}
\begin{tikzpicture}
    \def\gIpIIIs{0.6}
    \begin{scope}
        \gIpIIIgrid{owo,ooo}
        \draw[gtree] (0,0) -- (0,-0.6);
        \draw[gtree] (0,-0.6) -- (0.6,-0.6);
        \draw[gtree] (0.6,-0.6) -- (1.2,-0.6);
        \draw[gtree] (1.2,0) -- (1.2,-0.6);
        \node[font=\scriptsize\bfseries, align=center] at (0.6,-1.4) {unite with the land\\below and to the right};
    \end{scope}
    \draw[-{Stealth[length=3mm]}, line width=1.2pt, primaryblue] (2.0,-0.3) -- (3.0,-0.3);
    \begin{scope}[xshift=3.6cm, yshift=0.3cm]
        \node[gvdst, minimum size=7mm, font=\scriptsize] (rt) at (1.2,0) {};
        \node[gv, minimum size=7mm, font=\scriptsize] (a) at (0,-1.2) {};
        \node[gv, minimum size=7mm, font=\scriptsize] (b) at (0.8,-1.2) {};
        \node[gv, minimum size=7mm, font=\scriptsize] (c) at (1.6,-1.2) {};
        \node[gv, minimum size=7mm, font=\scriptsize] (d) at (2.4,-1.2) {};
        \foreach \x in {a,b,c,d} \draw[gde] (\x) -- (rt);
        \node[gdist, anchor=west] at ([xshift=2pt]rt.east) {root: sz $=5$};
    \end{scope}
    \node[note, anchor=west, text width=5.2cm, align=left] at (7.6,-0.4) {One DSU tree per island; its root stores the island's size. The exact shape depends on the union order (path compression flattens it), the size does not.};
\end{tikzpicture}
\figcap{DSU view of a 5-cell island: no traversal at all, only ``merge two sets'' operations, each adding the sizes.}
\end{diagrambox}
All four versions agree with each other on $20\,000$ random grids up to $7\times7$ (we ran them).
\end{alternativebox}

\begin{tricksbox}
\begin{itemize}[leftmargin=*]
    \item \textbf{Bigger grids:} with $RC$ up to $10^6$ replace the recursion by the explicit stack of Problem 1.3.1; the area is simply the number of pops.
    \item \textbf{Keep the input:} if the caller still needs \texttt{grid}, sink with the value $2$ and turn all $2$s back into $1$ at the end, or use \texttt{seen} as in the BFS version.
    \item \textbf{Early stop:} once \texttt{best} is at least the number of land cells not yet sunk, no later island can beat it --- rarely needed, but a cheap bound when the grid is huge and mostly one island.
\end{itemize}
\end{tricksbox}

\subsubsection*{5. Trace}

\begin{tracebox}[{Example 1 --- the scan and the fill of island E}]
\begin{center}
\small
\begin{tabular}{@{}clcc@{}}
\toprule
\textbf{island} & \textbf{scan meets} & \textbf{\texttt{sink} returns} & \textbf{\texttt{best} after} \\
\midrule
A & $(0,2)$ & 1 & 1 \\
B & $(0,7)$ & 4 & 4 \\
C & $(2,1)$ & 4 & 4 \\
D & $(2,4)$ & 5 & 5 \\
E & $(3,8)$ & 6 & \textbf{6} \\
F & $(6,7)$ & 5 & 6 \\
\bottomrule
\end{tabular}
\end{center}

\begin{center}
\begin{tikzpicture}
    \def\gIpIIIs{0.5}
    \begin{scope}
        \gIpIIIgrid{wwkwwwwkwwwww,wwwwwwwkkkwww,wkkwkwwwwwwww,wkwwkkwwxwoww,wkwwkkwwoooww,wwwwwwwwwwoww,wwwwwwwooowww,wwwwwwwoowwww}
        \node[font=\scriptsize\bfseries, align=center] at (3.0,-4.3) {the scan reaches $(3,8)$ (orange); A--D are\\already sunk (light blue = now water)};
    \end{scope}
    \begin{scope}[xshift=8.6cm]
        \node[gvsrc, minimum size=9mm, font=\tiny] (s1) at (0,0) {3,8};
        \node[gv, minimum size=9mm, font=\tiny] (s2) at (0,-1.15) {4,8};
        \node[gv, minimum size=9mm, font=\tiny] (s3) at (0,-2.3) {4,9};
        \node[gv, minimum size=9mm, font=\tiny] (s4) at (0,-3.45) {4,10};
        \node[gv, minimum size=9mm, font=\tiny] (s5) at (-1.0,-4.6) {3,10};
        \node[gv, minimum size=9mm, font=\tiny] (s6) at (1.0,-4.6) {5,10};
        \draw[gtree] (s1) -- node[lbl, left]{down} (s2);
        \draw[gtree] (s2) -- node[lbl, left]{right} (s3);
        \draw[gtree] (s3) -- node[lbl, left]{right} (s4);
        \draw[gtree] (s4) -- node[lbl, left]{up} (s5);
        \draw[gtree] (s4) -- node[lbl, right]{down} (s6);
        \foreach \n/\v in {s1/6, s2/5, s3/4, s4/3, s5/1, s6/1}
            \node[gdist, anchor=west] at ([xshift=1pt]\n.east) {$\to$ \v};
        \node[font=\scriptsize\bfseries, align=center] at (0.2,-5.5) {recursion tree of \texttt{sink(3,8)}};
    \end{scope}
\end{tikzpicture}
\end{center}
The call \texttt{sink(3,8)} goes down to $(4,8)$, right twice to $(4,10)$, and there splits: up to $(3,10)$ and down to $(5,10)$. The leaves return 1; on the way back every call adds 1 for itself: $1+1+1 = 3$, then 4, 5 and finally \textbf{6}. The maximum depth is 5 frames. The diagonal neighbour $(6,9)$ of $(5,10)$ is never called --- only side neighbours are --- so island F is measured separately (5) and the answer stays \textbf{6}.
\end{tracebox}

\subsubsection*{6. Test suite}

\begin{testbox}
\centering\small
\begin{tabular}{@{}lll@{}}
\toprule
\textbf{grid} & \textbf{Expected} & \textbf{Why} \\
\midrule
Example 1 & 6 & a diagonal contact must not merge E and F \\
\texttt{[[0,0,0,0,0,0,0,0]]} & 0 & no island \\
\texttt{[[1]]} & 1 & single cell \\
\texttt{[[1,1],[1,1]]} & 4 & the mark-on-pop bug returns 5 here \\
\texttt{[[1,0,1],[0,1,0],[1,0,1]]} & 1 & five islands touching at corners only \\
\texttt{[[1,1,0],[1,1,0],[0,0,1]]} & 4 & a square and a lone cell touching at a corner \\
$50\times50$ all \texttt{1} & 2500 & largest island, deepest recursion \\
$50\times50$ checkerboard & 1 & 1250 islands of area 1 \\
\bottomrule
\end{tabular}
\end{testbox}

\subsubsection*{7. Common mistakes}

\begin{warningbox}[Mistake 1 --- sinking after the recursive calls]
\begin{lstlisting}[style=snippet]
int a = 1 + sink(g, r - 1, c) + sink(g, r + 1, c) + ...;
g[r][c] = 0;                         // WRONG: too late, the neighbours call back into (r,c)
\end{lstlisting}
The neighbour's call still sees $(r,c)$ as land and recurses back: infinite recursion. Sink the cell \emph{first}.
\end{warningbox}

\begin{warningbox}[Mistake 2 --- forgetting the ``no island'' case]
Initialise \texttt{best = 0}, not with \texttt{INT\_MIN} or with the first island's area: the second example has no land and expects 0.
\end{warningbox}

\begin{warningbox}[Mistake 3 --- counting diagonal neighbours]
With 8 directions Example 1 returns 11 (E and F merge). The statement says ``horizontal or vertical''.
\end{warningbox}

\begin{warningbox}[Mistake 4 --- a global counter that is not reset]
A popular style keeps \texttt{int area} as a class member and increments it inside the fill. Forgetting \texttt{area = 0} before each new island makes every island ``inherit'' the previous ones: the answer becomes the total amount of land. Returning the area from the function, as above, makes this mistake impossible.
\end{warningbox}

\subsubsection*{8. Takeaways}

\begin{itemize}[leftmargin=*]
    \item Size of a component = number of cells a fill marks; a recursive fill can \emph{return} it as $1 + \sum$ children.
    \item Sink a cell \emph{before} recursing; check bounds before reading.
    \item Small grids ($\le 2500$ cells): recursion is fine. Large grids: explicit stack.
    \item Count, size, sum, maximum --- all are ``scan + fill + one number per fill''.
\end{itemize}

\begin{variantsbox}
\variant{Flip one water cell}{\varin{https://leetcode.com/problems/making-a-large-island/}}
  {you may turn at most one $0$ into $1$; maximise the largest island.}
  {label every island with an id and store its area; a water cell joins its distinct neighbouring islands.}
  {for each $0$: $1 + $ the areas of its \emph{distinct} neighbouring ids; answer = max with the largest island (Exercise E3.5).}
\variant{Fish instead of area}{\exmedium\ \lc{Medium}\ \varlink{Maximum Number of Fish in a Grid}{LeetCode 2658}{https://leetcode.com/problems/maximum-number-of-fish-in-a-grid/}}
  {water cells hold fish; collect the most fish in one connected water region.}
  {vertices = cells with a positive value; same 4-neighbour edges.}
  {the fill returns the sum of the values instead of the number of cells.}
\variant{Perimeter instead of area}{\varin{https://leetcode.com/problems/island-perimeter/}}
  {measure the boundary of the (single) island.}
  {count edges between a land cell and water or the outside, instead of vertices.}
  {$4\cdot\#\text{land} - 2\cdot\#\text{land--land edges}$, no traversal needed (Exercise E3.1).}
\variant{Histogram of island sizes}{\varnone}
  {report how many islands have each area $1, 2, \dots$}
  {same components.}
  {\texttt{cnt[sink(r,c)]++} for every start of the scan; $\bigO{RC}$.}
\variant{Largest 8-connected island}{\varnone}
  {cells touching at a corner belong to the same island.}
  {edges also between diagonal neighbours.}
  {8 recursive calls instead of 4; Example 1 then returns 11.}
\end{variantsbox}
